%************************************************
\chapter{Summary}
\label{ch:conclusion} %
%************************************************

{\sffamily\slshape The text in this chapter is based on Section 6 of the published versions of the preprints at \href{https://arxiv.org/abs/2411.15900}{arXiv:2411.15900} and \href{https://arxiv.org/abs/2504.10609}{arXiv:2504.10609} and  Section 5 of the preprint at \href{https://arxiv.org/abs/2504.10609}{arXiv.2607.04160}.\\}

Thermodynamics emerged from efforts to analyze the operation of steam engines while maximizing their work.
In a nice analogy, the search for pathways may be phrased as viewing reaction networks as chemical engines in which we want to maximize the output of some target molecule(s).
Recognizing that thermodynamics, encapsulated by the Gibbs free energies of the reactions, play a pivotal role in determining the probability of a reaction, Chapter~\ref{ch:thermodynamic_model} focused on developing ways to integrate them into the search for pathways.

Exploring pathways with favorable energy profiles involves challenging navigation of a large and complex energy landscape where the driving thermodynamic forces are influenced by the concentration ratios, inflows of the source molecules, and outflows of the target molecules.
Chapter~\ref{ch:thermodynamic_model} attempted to address part of that challenge.
However, there might be several refinements possible to the model proposed in Chapter~\ref{ch:thermodynamic_model}, some of which we list here:
\begin{enumerate}
    \item The constraint enforcing that only reactions with negative free energies are allowed in the returned pathways (implemented through the constraint in Equation~\eqref{proposed_constraint3}) may be overly stringent.
    Under experimental conditions, the reactions observed to occur do not always align with those predicted to occur theoretically: the thermal decomposition of lithium carbonate is observed to occur in practice despite being predicted to be thermodynamically unfavorable using computed (or experimental) free energies for the reaction.
    However, reactions predicted to be thermodynamically favorable do not always occur in practice because they have a high energy barrier (like the allotropic conversion of diamond into graphite).
    Therefore, there should be some flexibility in delineating the boundary between those reactions that are considered likely and included in the pathway and those that are not, to bring the model closer to reality. 
    Such extensions of the current modeling should be fairly straightforward.
    \item The objective function minimizes the sum of the Gibbs free energies of the reactions in the pathway.
    However, this approach might favor a pathway with more extreme values for the free energies rather than one with more moderate free energy values.
    Therefore, it might be worthwhile to explore alternative objective functions, such as the min-max, in addition to the min-sum used here.
    \item The current formulation disregards the kinetics of individual reactions.
    However, the probability of a reaction occurring is determined to a substantial extent by the reaction barrier height, and not only by the concentration of the reactants and the products.
    As a result, our framework might include a reaction in the returned pathway with a favorable thermodynamic driving force but too high a reaction barrier, making it a slow and rate-limiting step for the entire pathway.
    On the other hand, coming up with a generic `kinetic oracle' with wide applicability to estimate the barrier heights for diverse reactions poses a significant challenge.
    The transition states determining the barrier heights are saddle points on the potential energy surface, unlike reactants or products that are minima and can be located by gradient descent optimization methods. 
\end{enumerate}

Hence, we do not claim that the model proposed in Chapter~\ref{ch:thermodynamic_model} to be the final in the quest to incorporate thermodynamics into the search for pathways in reaction networks.
Still, we believe that it constitutes a valuable step forward, and we end by summarizing its virtues.
We provide a systematic and thermodynamically informed approach to optimize the computational search for pathways in reaction networks.

Modeling pathways as integer hyperflows gives results that can be interpreted in mechanistic terms,
as each molecule is present in integral quantities.
Formulating the problem as an MILP problem and utilizing an MILP solver allows us to obtain solutions with reasonable computational resources and time.
It allows enumeration of pathways over the integral flow variables, facilitating a systematic exploration of pathways in the reaction network.
Returning several solutions to pathway queries is an asset, as the answer becomes more robust to possible errors in the calculations of chemical potentials of the molecules and to other artifacts of the simplification of reality inherent in any modeling.
Moreover, it offers alternatives to choose from if the optimal solution turns out to be unattainable under controlled industrial reactor or laboratory conditions. 

Finally, the application of the framework formulated in Chapter~\ref{ch:thermodynamic_model} to a reaction network, in Chapter~\ref{ch:casestudy1}, based on real-life chemistry acts as a proof-of-concept and illustrates the computational feasibility and analytic potential of our methods.

In Chapter~\ref{ch:refined_model}, we have described a methodology for kinetically informed exploration of pathways in reaction networks building on the model proposed in Chapter~\ref{ch:thermodynamic_model}. 
One part of the methodology is an automated workflow for annotating individual reactions with their energy barriers, computed using a diverse selection of tools---RDKit~\cite{rdkit_embedMolecule}, xTB~\cite{grimme_xtb}, ASE~\cite{ase-paper} and NeuralNEB~\cite{neuralneb}. The goal of this workflow is to minimize the laborious manual interventions and intensive expert knowledge often necessary when executing quantum chemical calculations to determine energy barriers. 
As demonstrated in Chapter~\ref{ch:casestudy2}, the workflow can be applied on networks of rather large scale, can work on the fly with generative network construction, and uses relatively moderate computational resources.
The method does not take into account the interactions of solvent molecules while calculating the energy barriers. This and the question of how well models based on machine learning generalize past their training sets may lead to reasonable reservations on the precision of the estimates of the energy barriers for some use cases. However, the proposed methodology offers the flexibility that this component can be substituted with alternatives, if desired.

As the other part of the proposed methodology, we formulated an ILP to query pathways in the annotated network and designed an objective function
for pathway evaluation. The query formalism is at the same time both precise and very flexible in what queries can be formulated. The objective function is made under the simplifying assumption that all molecules in the reaction network are present with unit concentrations at all times. Additionally, the costs assigned to pathways depend on the underlying reaction network considered, so comparisons across networks are not directly possible. However, the question of how to evaluate pathways is complex and the reader of course has the option to substitute a different objective function, based on other cost metrics. We do note, though, that our simplification and the formulation via ILP allows us to analyze much larger networks and using much less computational resources than alternatives like stochastic simulations or potential energy surface exploration.
We also note that by a progressive constraining of the ILP, we in the end can return a \emph{list} of high-scoring solutions, each representing structurally different pathways. Giving a list of good, structurally different solutions, rather than a single solution, helps mitigate disagreements about the details of pathway cost metrics.

Finally, the efficiency and utility of the method described in Chapter~\ref{ch:refined_model} was shown on a concrete reaction network in Chapter~\ref{ch:casestudy2}, studying pathways to pre-biotically significant molecules, such as glycine and 2-hydroxyethanoic acid. 

We have presented a framework for exploring reaction networks in Chapter~\ref{ch:stoc_explore}, motivated by chemical stochastic kinetics generating random walks on networks~\cite{randomwalk1,randomwalk2,randomwalk3} that differ from previous approaches. We do not precompute the network but rather discover regions of it through a process mimicking chemical stochastic dynamics. Instead of computing the relative probabilities for all possible reactions at each step, our method performs three modular sampling operations: (1) selecting a colliding molecular pair based on current molecular abundances, (2) selecting a reaction template, and (3) selecting a specific reaction instance generated by a graph transformation engine acting on that pair. In the example study presented in Chapter~\ref{ch:stoc_explore}, the reaction template probabilities are not intended to approximate physical rate laws, but rather to provide an unbiased sampling of dynamically reachable regions of a reaction graph defined on an atomistic level. This divide-and-conquer strategy avoids enumerating the full list of transitions out of the current state at every step and allows the dynamics to generate previously unvisited regions of chemical space on demand. In this sense, the method performs a biased random walk on an implicitly defined molecular network, constructing only the locally relevant subnetworks during the course of the exploration. 

Some key features of the framework proposed in the current work include the simultaneous use of (1) stochastic simulation-styled dynamics operating on molecular graphs (2) applying rule-based bond rearrangements on an atomic level using reaction templates to generate a relevant part of the molecular network on the fly and (3) the support of open systems with inflow and outflow on the system. 
It borrows from and straddles the region between the Gillespie algorithm for stochastic simulations (which is kinetically accurate but abstract), 
generative chemistry using graph grammar (which is chemically accurate but deterministic) and kinetic Monte Carlo simulations (which are stochastic but non-generative and closed). 
The simulation can be steered to explore the region of the molecular space under interest--- its parameters might also be biased to explore events which rarely occur under realistic conditions but are nevertheless important.
The present method provides a lightweight, highly scalable way to navigate regions of chemical space where kinetic parameters are unknown or where the objective is discovery rather than quantitative prediction.
It enables stochastic, network-free exploration of large, generative molecular spaces under minimal kinetic assumptions.
Together, with stochastic simulation schemes like the Gillespie algorithm or kinetic Monte Carlo, it forms two complementary probes to molecular networks: one for exact stochastic dynamics, one for graph-based exploration.

In summary, this work contributes to field of stochastical exploring large reaction networks and computationally assisted analysis and design of pathways in reaction networks represented as directed hypergraphs.

